\chapter{Внутрішня будова та життєвий цикл шкідливого коду}
\chaptershort{Внутрішня будова шкідливого коду}
\label{chap:ch05}

Ефективна протидія злоякісному коду неможлива без розуміння того, як він улаштований і як виконується. Аналітик, що досліджує зразок, розробник правил виявлення, спеціаліст із реагування --- усі спираються на знання форматів виконуваних файлів, механізмів завантаження в пам'ять, способів закріплення в системі й приховування присутності. Це знання має суто захисний характер: воно дозволяє визначати, де шукати ознаки компрометації, які артефакти свідчать про зараження, як відрізнити легітимну поведінку від злоякісної. Розділ розглядає внутрішню будову й життєвий цикл шкідливого коду з позиції аналітика-захисника: формати виконуваних файлів, механізми завантаження й виконання, техніки ін'єкції коду та способи їх виявлення, механізми персистентності, підвищення привілеїв, приховування (руткіти) й модульну архітектуру сучасного ЗПК. Виклад орієнтовано на \emph{розпізнавання} й \emph{аналіз}, а не на створення коду. Інструменти статичного й динамічного аналізу, що спираються на цей матеріал, розглянуто в розділах~\ref{chap:ch06} і~\ref{chap:ch07}.

\section{Формати виконуваних файлів}\label{sec:ch05-formats}

\subsection{Структура формату PE}

Основним форматом виконуваних файлів Windows є PE (Portable Executable)~\cite{ch05:pecoff}. Знання його структури є фундаментом статичного аналізу, оскільки більшість ЗПК для найпоширенішої платформи має цей формат. PE-файл складається з (рис.~\thechapter.1): заголовка DOS (історична сумісність), сигнатури PE, файлового заголовка COFF (архітектура, кількість секцій, час компіляції), опційного заголовка (точка входу, образ бази, вирівнювання, характеристики), таблиці секцій і самих секцій. Типові секції наведено в табл.~\thechapter.1.

\tabhead{Таблиця~\thechapter.1 --- Типові секції PE-файлу та їхнє значення для аналізу}

\begin{center}\footnotesize
\begin{tabularx}{\linewidth}{L{16mm} Y}
\toprule
Секція & Вміст і значення для аналізу \\
\midrule
\texttt{.text} & виконуваний код; аномально висока ентропія свідчить про пакування \\
\texttt{.data} & ініціалізовані дані; рядки, конфігурація \\
\texttt{.rdata} & константи, таблиця імпортів (підказує функціональність) \\
\texttt{.rsrc} & ресурси; часто містять вкладені навантаження \\
\texttt{.reloc} & релокації; відсутність може вказувати на особливості завантаження \\
нетипові & нестандартні імена секцій --- ознака пакувальника \\
\bottomrule
\end{tabularx}
\end{center}

Для аналітика найцінніша інформація --- \term{таблиця імпортів} (список функцій, які програма викликає з системних бібліотек): вона підказує функціональність ще до виконання. Наприклад, імпорт функцій роботи з криптографією й файловою системою разом характерний для вимагачів; функцій ін'єкції й маніпуляції процесами --- для імплантів; мережевих функцій і функцій реєстру --- для бекдорів. Набір імпортів утворює характерний відбиток, що використовується для ідентифікації сімейств через \term{imphash} --- геш упорядкованого списку імпортів~\cite{ch05:imphash}: зразки одного сімейства часто мають однаковий imphash, що дозволяє їх кластеризувати (розділ~\ref{chap:ch06}).

\subsection{Інші формати}

Для Unix-подібних систем (Linux, зокрема сервери й IoT, розділ~\ref{chap:ch14}) основним є формат ELF~\cite{ch05:elf} з аналогічною логікою: заголовок, програмні заголовки (сегменти для завантаження), заголовки секцій, таблиці символів і релокацій. Для macOS --- Mach-O. Крім нативних форматів, значна частина сучасного ЗПК використовує \emph{проміжні} й \emph{інтерпретовані} форми: байт-код .NET і Java, скрипти (PowerShell, JavaScript, VBA-макроси в документах), що потребують інших підходів до аналізу (розділ~\ref{chap:ch13}). Документи офісних форматів із макросами та експлойтами є поширеним носієм для доставки (розділ~\ref{chap:ch01}), тому їх структура (OLE, OOXML) також є об'єктом аналізу.

\begin{tikzfig}
\begin{tikzpicture}[font=\scriptsize,
  b/.style={draw,minimum width=34mm,minimum height=5.5mm,align=center},
  >={Stealth[length=1.5mm]}]
\node[b] (dos) at (0,0) {Заголовок DOS + заглушка};
\node[b,below=0.6mm of dos] (pe) {Сигнатура PE + заголовок COFF};
\node[b,below=0.6mm of pe] (opt) {Опційний заголовок (точка входу)};
\node[b,below=0.6mm of opt] (sect) {Таблиця секцій};
\node[b,below=0.6mm of sect,fill=black!6] (t) {\texttt{.text} (код)};
\node[b,below=0.6mm of t,fill=black!6] (d) {\texttt{.data}, \texttt{.rdata} (дані, імпорти)};
\node[b,below=0.6mm of d,fill=black!6] (r) {\texttt{.rsrc} (ресурси)};
\node[anchor=west,font=\tiny,align=left] at (2.1,0) {метадані:};
\node[anchor=west,font=\tiny,align=left] at (2.1,-0.95) {архітектура,};
\node[anchor=west,font=\tiny,align=left] at (2.1,-1.3) {час компіляції,};
\node[anchor=west,font=\tiny,align=left] at (2.1,-1.65) {точка входу};
\end{tikzpicture}
\figcap{Рис.~\thechapter.1 --- Структура PE-файлу (за \cite{ch05:pecoff})}
\end{tikzfig}

\section{Завантаження та виконання}\label{sec:ch05-loading}

Розуміння процесу завантаження є ключем до виявлення аномалій. Легітимний виконуваний файл завантажується системним завантажувачем: образ відображається в пам'ять процесу за адресою бази, виконуються релокації (коригування адрес), розв'язуються імпорти (адреси функцій заповнюються в таблиці імпортів), налаштовуються права доступу до сторінок пам'яті (код --- для виконання, дані --- без виконання), і керування передається точці входу. Будь-яке відхилення від цього штатного шляху є потенційним індикатором злоякісності.

Шкідливий код часто завантажується в обхід штатного механізму, щоб уникнути залишення файлу на диску й виявлення. Типові аномалії, що є ознаками для виявлення: ділянки пам'яті з правами «запис і виконання» одночасно (RWX), характерні для коду, згенерованого в пам'яті; виконання коду з ділянок, не підкріплених файлом на диску (приватна, а не відображена пам'ять); невідповідність між секціями на диску й у пам'яті; розв'язання імпортів вручну (через пошук функцій за гешами імен замість таблиці імпортів) --- техніка, що приховує справжню функціональність від статичного аналізу. Сучасні засоби захисту кінцевих точок (розділ~\ref{chap:ch10}) моніторять саме ці аномалії завантаження, оскільки вони стійкіші до змін коду, ніж сигнатури файлів.

Окремим важливим механізмом є \emph{динамічне розв'язання} системних функцій. Легітимна програма зазвичай оголошує свої залежності в таблиці імпортів, що робить її функціональність прозорою для статичного аналізу. Шкідливий код натомість часто завантажує бібліотеки й отримує адреси функцій під час виконання за іменами або їх гешами, тому таблиця імпортів виглядає «порожньою» або оманливо безневинною. Для аналітика це подвійний сигнал: з одного боку, бідна таблиця імпортів за наявності складного коду сама собою підозріла; з іншого --- щоб відновити справжню функціональність, потрібен динамічний аналіз (розділ~\ref{chap:ch07}) або емуляція. Розуміння цього механізму пояснює, чому суто статичний аналіз таблиці імпортів недостатній і чому сучасні засоби поєднують статичні й поведінкові ознаки. Ще однією поширеною технікою є виконання \term{шелкоду} --- позиційно-незалежного коду без структури PE, що ускладнює його розпізнавання як виконуваного; виявлення шелкоду спирається на аналіз вмісту пам'яті й поведінки, а не файлової структури.

\section{Техніки виконання й ін'єкції коду}\label{sec:ch05-injection}

Для прихованого виконання й доступу до ресурсів інших процесів ЗПК застосовує \term{ін'єкцію коду} --- виконання власного коду в контексті іншого (зазвичай легітимного) процесу~\cite{ch05:barnea,ch05:klein,ch05:attck-t1055}. Мета --- маскування (злоякісна активність виглядає як дія довіреної програми), обхід засобів контролю застосунків, доступ до даних цільового процесу (наприклад, браузера для крадіжки сесій). З позиції захисту важливо знати \emph{сліди}, які залишають ці техніки, узагальнені в табл.~\thechapter.2 (виклад описує принцип і ознаки виявлення, а не реалізацію).

\tabhead{Таблиця~\thechapter.2 --- Техніки ін'єкції коду та ознаки їх виявлення}

\begin{center}\footnotesize
\begin{tabularx}{\linewidth}{L{28mm} Y}
\toprule
Техніка (принцип) & Характерні ознаки для виявлення \\
\midrule
Ін'єкція бібліотеки в процес & завантаження бібліотеки з незвичного шляху в чужий процес \\
Запис і виконання в чужому процесі & виділення пам'яті RWX у чужому процесі, створення віддаленого потоку \\
Підміна образу процесу (hollowing) & невідповідність образу в пам'яті файлу на диску; призупинений запуск \\
Відображення коду без файлу & виконання з приватної пам'яті, ручне розв'язання імпортів \\
Зловживання механізмами ОС & аномальне використання легітимних системних механізмів розширення \\
\bottomrule
\end{tabularx}
\end{center}

Принципова складність виявлення ін'єкції полягає в тому, що майже кожна її елементарна операція (виділення пам'яті, запис у процес, створення потоку) має легітимні застосування: налагоджувачі, засоби захисту, системи віртуалізації, менеджери вікон законно виконують подібні дії. Тому виявлення спирається не на сам факт операції, а на \emph{нетиповий контекст}: хто саме, в який процес, за яких обставин і в якій послідовності її виконує. Наприклад, текстовий редактор, що виділяє RWX-пам'ять у процесі браузера й створює там потік, --- очевидна аномалія, тоді як та сама дія від налагоджувача щодо програми, яку він налагоджує, --- норма. Це робить виявлення ін'єкції задачею кореляції послідовності подій, а не розпізнавання окремої дії, і пояснює, чому ефективні детектори потребують повної телеметрії процесів (розділ~\ref{chap:ch10}).

Для кількісної пріоритизації підозрілих подій ін'єкції застосовують \emph{скорингові} евристики. Нехай подія характеризується набором бінарних ознак $f_1,\dots,f_n$ (RWX-пам'ять, крос-процесність, невідповідність диску, ручні імпорти тощо) з вагами $w_i$. Тоді оцінка підозрілості

\begin{equation}
S = \sum_{i=1}^{n} w_i\, f_i,
\label{eq:ch05-score}
\end{equation}

а тривога генерується за $S \ge \theta$. Ваги калібрують так, щоб мінімізувати хибні тривоги на легітимній активності (деякі програми легітимно використовують подібні механізми), що знову приводить до задачі базової частоти (розділ~\ref{chap:ch03}): жодна окрема ознака не є достатньою, і рішення ухвалюється за сукупністю. Саме такий підхід реалізують поведінкові детектори EDR (розділ~\ref{chap:ch10}).

\section{Персистентність}\label{sec:ch05-persistence}

Щоб пережити перезавантаження й зберегти доступ, ЗПК забезпечує \term{персистентність} --- автоматичне відновлення виконання~\cite{ch05:attck-persist,ch05:russinovich}. Оскільки механізмів персистентності в сучасних ОС багато, їх систематичний моніторинг є ключовим елементом виявлення. Основні класи: автозапуск через ключі реєстру й теки автозавантаження; реєстрація як системної служби; заплановані завдання; підміна чи «викрадення» порядку завантаження легітимних бібліотек; модифікація механізмів запуску застосунків; на найглибшому рівні --- буткіти, що закріплюються в завантажувачі чи прошивці (розглянуто далі).

Ключова властивість для виявлення: персистентність \emph{залишає стійкий артефакт} у системі (запис у реєстрі, файл, завдання, службу), який можна виявити інвентаризацією точок автозапуску й порівнянням із відомим чистим станом. Тому один із найефективніших методів виявлення --- регулярний аудит усіх точок автозапуску (існують спеціалізовані утиліти, що перелічують десятки їх категорій) з перевіркою підписів і репутації. Це приклад загального принципу: кожна дія зловмисника, що забезпечує стійкість, створює й стійкий слід, який стає індикатором компрометації (розділ~\ref{chap:ch11}). Для безфайлового коду (розділ~\ref{chap:ch13}) персистентність реалізується в реєстрі чи через заплановані завдання без файлу на диску, що ускладнює виявлення й вимагає поведінкового моніторингу.

Вибір механізму персистентності зловмисником залежить від компромісу між \emph{надійністю}, \emph{скритністю} й потрібними \emph{привілеями}. Прості механізми (теки автозавантаження, ключі автозапуску реєстру) не потребують адміністративних прав, але легко виявляються; складніші (служби, драйвери, буткіти) потребують вищих привілеїв, зате стійкіші й менш помітні. Державні угруповання (розділ~\ref{chap:ch13}) надають перевагу скритним механізмам і нерідко зловживають легітимними функціями розширення системи, щоб злитися з нормальною активністю. Для захисника це означає, що виявлення персистентності не зводиться до перевірки кількох «популярних» місць: потрібне систематичне покриття всіх класів механізмів і, що важливіше, \emph{базова лінія} нормального стану системи, відносно якої виявляються відхилення. Побудова й підтримка такої базової лінії (відомі легітимні автозапуски, підписані драйвери, звичайні заплановані завдання) є практичною основою як для ручного полювання на загрози, так і для автоматизованого виявлення (розділ~\ref{chap:ch11}). Саме відсутність базової лінії робить виявлення персистентності в реальних інфраструктурах складнішим, ніж здається в теорії.

\section{Підвищення привілеїв}\label{sec:ch05-privesc}

Для повного контролю над системою ЗПК прагне підвищити привілеї від рівня користувача до адміністратора чи ядра. Шляхи: експлуатація вразливостей ядра чи системних служб; зловживання некоректними конфігураціями (надмірні права, вразливі служби); обхід механізмів контролю доступу; крадіжка й повторне використання токенів доступу. З позиції захисту підвищення привілеїв протидіє принцип найменших привілеїв (мінімізація прав), своєчасне усунення вразливостей (розділ~\ref{chap:ch10}), моніторинг аномального використання привілеїв і захист облікових даних (запобігання їх вилученню з пам'яті). Виявлення спроб підвищення привілеїв є важливим сигналом, оскільки вказує на перехід атаки від початкового доступу до закріплення й розширення контролю (розділ~\ref{chap:ch04}).

\section{Приховування: руткіти}\label{sec:ch05-rootkits}

\term{Руткіт} --- це компонент, що приховує присутність ЗПК у системі, модифікуючи її так, щоб злоякісні файли, процеси, мережеві з'єднання й записи реєстру не відображалися звичайними засобами~\cite{ch05:hoglund,ch05:blunden}. За рівнем дії розрізняють руткіти режиму користувача (перехоплення викликів у застосунках) і режиму ядра (модифікація структур і функцій ядра --- найнебезпечніші, бо діють на рівні, що контролює засоби захисту). Концептуальний приклад техніки ядра --- пряма маніпуляція об'єктами ядра (приховування процесу вилученням його з переліку), що робить процес невидимим для стандартних засобів переліку, але лишає сліди в інших структурах.

Протидія руткітам ґрунтується на кількох принципах. \emph{Перехресна перевірка} (cross-view): порівняння результатів, отриманих різними способами. Формально, якщо $\mathcal{V}_H$ --- множина об'єктів (процесів, файлів, з'єднань), отримана високорівневим інтерфейсом, а $\mathcal{V}_L$ --- незалежним низькорівневим, то приховані руткітом об'єкти виявляються як різниця

\begin{equation}
\mathcal{H} = \mathcal{V}_L \setminus \mathcal{V}_H,
\label{eq:ch05-crossview}
\end{equation}

тобто об'єкти, видимі «знизу», але приховані «згори». Непорожня $\mathcal{H}$ є сильним індикатором активного приховування. Ефективність методу залежить від того, наскільки $\mathcal{V}_L$ справді незалежний від руткіта: якщо руткіт діє на рівні, нижчому за обидва інтерфейси, різниця буде порожньою, тому потрібен аналіз із максимально довіреного рівня. \emph{Аналіз із довіреного рівня}: дослідження системи з рівня, нижчого за руткіт (завантаження з чистого носія, аналіз образу пам'яті, апаратна підтримка). \emph{Захист цілісності ядра}: сучасні ОС застосовують обов'язковий підпис драйверів, захист цілісності коду ядра, віртуалізаційну ізоляцію критичних структур, що суттєво ускладнює руткіти режиму ядра. \emph{Довірене завантаження}: механізми Secure Boot і вимірюваного завантаження протидіють буткітам~\cite{ch05:matrosov}, перевіряючи цілісність ланцюга завантаження від апаратного кореня довіри. Попри це, найскладніші загрози (зокрема державні, розділ~\ref{chap:ch13}) досягали рівня прошивки, тому аналіз пам'яті й цілісності залишається важливим напрямом криміналістики (розділ~\ref{chap:ch16}).

\section{Модульна архітектура сучасного ЗПК}\label{sec:ch05-modular}

Сучасний злоякісний код рідко є монолітним. Типова архітектура складного ЗПК (рис.~\thechapter.2) є \emph{модульною}: невеликий початковий компонент (dropper/loader) забезпечує доставку й закріплення, завантажує з інфраструктури зловмисника основне ядро (бекдор із C2), а воно за командою підвантажує спеціалізовані модулі (крадіжка даних, рух мережею, шифрування). Така архітектура дає зловмиснику гнучкість (оновлення й заміна модулів), а захиснику створює складність (на диску може не бути повного функціоналу, частина завантажується лише в пам'ять за командою). Водночас модульність створює й можливості для виявлення: кожен модуль має взаємодіяти з іншими, завантажуватися й спілкуватися з C2, залишаючи поведінкові сліди на різних етапах ланцюга ураження (розділ~\ref{chap:ch04}).

\begin{tikzfig}
\begin{tikzpicture}[font=\scriptsize,
  b/.style={draw,rounded corners=2pt,minimum height=7mm,text width=22mm,align=center},
  m/.style={draw,rounded corners=2pt,minimum height=6mm,text width=18mm,align=center,font=\tiny},
  >={Stealth[length=1.5mm]}]
\node[b] (drop) at (0,0) {Завантажувач (dropper/loader)};
\node[b] (core) at (3.6,0) {Ядро: бекдор + C2};
\node[m] (m1) at (7,0.9) {модуль крадіжки даних};
\node[m] (m2) at (7,0) {модуль руху мережею};
\node[m] (m3) at (7,-0.9) {модуль впливу};
\draw[->] (drop) -- node[above,font=\tiny] {завантажує} (core);
\draw[->] (core) -- (m1); \draw[->] (core) -- (m2); \draw[->] (core) -- (m3);
\node[font=\tiny] at (3.6,-1.1) {зв'язок з інфраструктурою зловмисника};
\end{tikzpicture}
\figcap{Рис.~\thechapter.2 --- Модульна архітектура сучасного складного злоякісного коду}
\end{tikzfig}

\section{Кейс: аналіз будови типового імпланта}\label{sec:ch05-case}

Розглянемо послідовність аналізу будови зразка з позиції захисника (без виконання, статично --- розділ~\ref{chap:ch06}).

\emph{Крок 1. Формат і метадані.} Визначаємо тип файлу (PE/ELF), архітектуру, час компіляції (може бути підробленим), наявність підпису. Обчислюємо криптографічний геш (для порівняння з базами) та imphash (для пошуку сімейства~\cite{ch05:imphash}).

\emph{Крок 2. Ентропія й пакування.} Аналізуємо ентропію секцій~\eqref{eq:ch03-entropy}: висока ентропія \texttt{.text} і нетипові імена секцій свідчать про пакування (розділ~\ref{chap:ch08}), що потребує розпакування перед глибоким аналізом.

\emph{Крок 3. Імпорти й рядки.} Таблиця імпортів підказує функціональність (мережа, криптографія, ін'єкція). Рядки (URL, шляхи, повідомлення) дають індикатори, хоча часто зашифровані.

\emph{Крок 4. Персистентність і виконання.} Виявляємо, які механізми закріплення й виконання використовує зразок (за імпортами, рядками, а точніше --- динамічним аналізом), щоб визначити артефакти для пошуку в заражених системах.

\emph{Крок 5. Оцінка підозрілості.} За сукупністю ознак (пакування, підозрілі імпорти, аномальні секції) обчислюємо оцінку~\eqref{eq:ch05-score} й ухвалюємо рішення про класифікацію й пріоритет.

\emph{Результат.} Будова зразка дає: сімейство (за imphash і подібністю), ймовірну функціональність (за імпортами), індикатори компрометації (геші, механізми персистентності, артефакти), пріоритет реагування. Ці результати живлять правила виявлення (розділ~\ref{chap:ch06}) і реагування (розділ~\ref{chap:ch16}). Аналіз будови реальних імплантів APT-угруповань проти України систематизовано у вітчизняних дослідженнях~\cite{ch05:apt-analysis}, а засоби автоматизованого аналізу ПЗ на небезпечні конструкції --- у~\cite{ch05:windows-vuln}. Кейс ілюструє, що знання внутрішньої будови перетворює непрозорий бінарний файл на набір конкретних індикаторів для захисту.

\section*{Висновки до розділу}
\addcontentsline{toc}{section}{Висновки до розділу}

Розуміння внутрішньої будови й життєвого циклу шкідливого коду є фундаментом його аналізу й виявлення. Формати виконуваних файлів (PE, ELF) мають структуру, що розкриває функціональність ще до виконання: таблиця імпортів підказує можливості, а imphash дозволяє кластеризувати сімейства; аномальна ентропія й нетипові секції свідчать про пакування. Процес завантаження має штатний шлях, будь-яке відхилення від якого (RWX-пам'ять, виконання без файлу, ручні імпорти) є стійким індикатором злоякісності, стійкішим за сигнатури файлів. Техніки ін'єкції коду забезпечують маскування й доступ до чужих процесів, але залишають характерні сліди, які виявляють поведінкові детектори за сукупністю ознак~\eqref{eq:ch05-score}. Персистентність неминуче створює стійкий артефакт у системі, що робить аудит точок автозапуску ефективним методом виявлення. Руткіти приховують присутність модифікацією системи, і протидія їм ґрунтується на перехресній перевірці, аналізі з довіреного рівня, захисті цілісності ядра й довіреному завантаженні. Сучасний ЗПК є модульним, що дає зловмиснику гнучкість, а захиснику --- додаткові поведінкові сліди на різних етапах. Головний висновок: кожна дія зловмисника, спрямована на виконання, закріплення чи приховування, створює артефакт або поведінкову ознаку, і саме знання будови перетворює ці ознаки на конкретні індикатори для захисту.

\selfcheck
\begin{enumerate}
\item Опишіть структуру PE-файлу (рис.~\thechapter.1). Яку інформацію дає таблиця імпортів і що таке imphash?
\item Чому висока ентропія секції \texttt{.text} і нетипові імена секцій є ознаками пакування (табл.~\thechapter.1)?
\item Опишіть штатний процес завантаження виконуваного файлу. Які відхилення є індикаторами злоякісності?
\item Чому ознаки аномалій завантаження (RWX, виконання без файлу) стійкіші за сигнатури файлів?
\item З якою метою ЗПК застосовує ін'єкцію коду? Які сліди вона залишає (табл.~\thechapter.2)?
\item Поясніть скорингову евристику~\eqref{eq:ch05-score}. Чому жодна окрема ознака ін'єкції не є достатньою для тривоги?
\item Назвіть класи механізмів персистентності. Чому персистентність завжди залишає стійкий артефакт і як це використовують для виявлення?
\item Чому аудит точок автозапуску є ефективним методом виявлення? Як ускладнює виявлення безфайлова персистентність?
\item Якими шляхами ЗПК підвищує привілеї? Які заходи протидіють підвищенню привілеїв?
\item Що таке руткіт? Порівняйте руткіти режиму користувача й ядра за рівнем небезпеки.
\item Опишіть принципи протидії руткітам. Як довірене завантаження протидіє буткітам?
\item Опишіть модульну архітектуру сучасного ЗПК (рис.~\thechapter.2). Які переваги вона дає зловмиснику й які можливості --- захиснику?
\item Опишіть послідовність аналізу будови зразка (кейс підрозділу~\ref{sec:ch05-case}). Які індикатори вона дає для захисту?
\item Дослідницьке: розробіть методику автоматичного виявлення аномалій завантаження (RWX-пам'ять, невідповідність диску) у телеметрії кінцевих точок і оцініть її частку хибних тривог з урахуванням легітимних випадків (розділ~\ref{chap:ch03}).
\item Дослідницьке: побудуйте систему кластеризації зразків за imphash і подібністю імпортів для автоматичної атрибуції сімейств та оцініть її стійкість до технік зміни імпортів (розділ~\ref{chap:ch08}).
\end{enumerate}
